\documentclass[main]{subfiles}
\usepackage[sols]{workbook}
\numbering{ws3-3}

\begin{document}

\ifSubfilesClassLoaded{%
  \chaptermark{\chapterthree}
}{}

\section{Autonomous First Order Differential Equations} \label{ws3-3}

\begin{framed}
\centerline{\textbf{Lesson Goals}}

You should be able to:
\begin{itemize}
\item Find the equilibrium solutions of a differential equation of the form $\frac{dy}{dt}=f(y)$
\item Use qualitative analysis to sketch representative solutions and specific solutions to a differential equation of the form $\frac{dy}{dt}=f(y)$
\item Use qualitative analysis to determine the long-term behavior of a quantity under specific initial conditions
\end{itemize}
\end{framed}

\begin{enumerate}
  \item 
  \note{This is just to emphasize again the idea that, if $\frac{dy}{dt}$
  depends only on $y$, then all the slopes in a horizontal line will be
  the same.}
  
 Here are three slope fields.

  \begin{tabular}{c c c}
  \includegraphics[width=\linewidth/3-\tabcolsep*2]{ws/images/slope-field-one-param-2-y-only.pdf}
  & 
  \includegraphics[width=\linewidth/3-\tabcolsep*2]{ws/images/slope-field-one-param-2-t-and-y}
  &
  \includegraphics[width=\linewidth/3-\tabcolsep*2]{ws/images/slope-field-one-param-2-t-only}
  \\
  (I) & (II) & (III)
  \end{tabular}
  
  For each of the following differential equations, find the picture above
  which shows its slope field.
  \probonly{\begin{multicols}{2}}
  \begin{enumerate}
    \item
    \begin{problem}
    $\displaystyle \frac{dy}{dt} = t^5 - 4t^3 + 4t - \frac{1}{2}$.
    \end{problem}

    \begin{solution}
    Here, $\frac{dy}{dt}$ depends only on $t$.  So, for instance, all
    points with $t = 1$ will have slope $\dfrac{1}{2}$.  What's
    important here is that all points on the vertical line $t = 1$
    will have the \textit{same} slope ($\frac{1}{2}$).  The only
    picture for which this is true is \fbox{(III)}. 
    \end{solution}

    \item
    \begin{problem}
    $\displaystyle \frac{dy}{dt} = y^5 - 4y^3 + 4y - \frac{1}{2}$.
    \end{problem}

    \begin{solution}
    Here, $\frac{dy}{dt}$ depends only on $y$.  So, for instance, all
    points with $y = 1$ will have slope $\dfrac{1}{2}$.  What's
    important here is that all points on the horizontal line $y = 1$
    will have the same slope ($\frac{1}{2}$).  The only picture for
    which this is true is \fbox{(I)}.
    \end{solution}

  \end{enumerate}
  \probonly{\end{multicols}}
  \vfill

%\pspace{\newpage}

\begin{problemonly}
\begin{framed}
\textbf{\centerline{Qualitative Analysis of Autonomous First-Order Differential Equations}}
To sketch solutions to a differential equation of the form $\displaystyle \frac{dy}{dx}=f(y)$:

\begin{enumerate}
  \item Find the equilibrium solutions by setting $\displaystyle \frac{dy}{dx}=f(y)=0$, and sketch these on the $xy$-plane.
  \item The equilibrium solutions divide the picture into horizontal
  strips.  In each strip, test one $y$-value to determine whether
  $\frac{dy}{dx}$ is positive or negative.  This tells you whether
  solutions in that strip increase or decrease. Draw
  a representative solution in each strip.
 \item If you want to draw the solution satisfying an initial condition, find the strip in which the initial condition appears. Then the solution sketch is just a horizontal shift of the representative solution in that strip. 
 \end{enumerate}
\end{framed}  
\end{problemonly}
  
  
\pspace{\newpage} 

\item 
\begin{problem}
  Suppose that a population of painted turtles in a pond can be modeled by the differential equation  
  $$\displaystyle \frac{dP}{dt} = \frac{1}{1000}P(300-P),$$ where $P(t)$ is the number of turtles and $t$ is measured in months from the present day. This is known as a \textit{logistic model} for population growth.
\end{problem}

\begin{enumerate}
  \item 
  \begin{problem}[2.2in]
    Use qualitative analysis to sketch a family of representative solutions to the differential equation. Draw the particular solution that  satisfies the initial condition $P(0)=50$ (i.e., the population right now is $50$ turtles). Please identify the stability/instability of any equilibrium solutions. What do you expect to happen to the turtle population over time?
  \end{problem}

  \begin{solution}
  The equilibrium solutions are $P=0$ (unstable) and $P=300$ (stable).
  The particular solution to the initial value problem with $P(0)=50$ is drawn in red.

  \includegraphics[width=2.5in]{ws/images/qualitative-eulers-method-turtles}
  
  In the long run, we should expect the turtle population to approach 300. 
  \end{solution}

\pspace{\vfill}

  \item 
  \begin{problem}[1in]
In what way(s) is this model more realistic than an exponential growth model? In what way(s) is this model unrealistic?
  \end{problem}
  
  \begin{solution}
  In an exponential growth model, the population increases exponentially without bound as time goes on. This is unrealistic for most species, as a population's growth is limited by environmental factors, such as space and availability of resources. So the logistic model is more realistic in the sense that it contains a \textit{carrying capacity} for the species, which the population will grow toward exponentially, but not exceed. 
  
  However, this logistic model also has its flaws. For example, if our initial condition is $P(0)=1$, our model predicts that the turtle population will grow towards 300 turtles. However, if we start with only one turtle, that turtle has no one to reproduce with! So realistically, the turtle population should eventually become 0.
  
  \end{solution}
   

\end{enumerate}

\pspace{\newpage}


  \item 
  Use qualitative analysis to sketch a family of representative solutions
  for each differential equation.  Identify the equilibrium solutions of
  each differential equation, and say whether they are stable or unstable.
  \probonly{\begin{multicols}{2}}
  \begin{enumerate}
    \item
    \begin{problem}[5in]
    $\displaystyle \frac{dy}{dt} = (y + 1) (4 - y)$
    \end{problem}

    \begin{solution}
    We start by finding the equilibrium solutions, where $\frac{dy}{dt}
    = (y + 1)(4 - y)$ is 0.  This happens when $y = -1$ or $y = 4$, so
    $y(t) = -1$ and $y(t) = 4$ are equilibrium solutions.

    Now, we need to figure out the sign of $\frac{dy}{dt}$ when $y <
    -1$, $-1 < y < 4$, and when $4 < y$.  To do this, we can test one
    point in each region, and we find:
    \begin{itemize}[topsep=0pt,itemsep=0pt]
    \item When $y < -1$, $\frac{dy}{dt} < 0$.  (For instance, plug in $y
    = -2$: $\frac{dy}{dt} = (y + 1) (4 - y)$ is equal to $(-1) (6) =
    -6$ here.)
    \item When $-1 < y < 4$, $\frac{dy}{dt} > 0$.
    \item When $4 < y$, $\frac{dy}{dt} < 0$.
  \end{itemize}
  We can put this information in a diagram:
  \begin{center}
  \includegraphics[width=2in]{ws/images/qa2-sketch}
  \end{center}
  The arrows indicate the general direction (up or down) of the
  solution curves in each strip.  Now, it is easy to draw a
  representative solution in each strip:
  \begin{center}
  \includegraphics[width=3in]{ws/images/qa2}
  \end{center}
  
Based on the solution curves above, the equilibrium solution $y(t)=4$ is \fbox{stable}, since solution curves above AND below the yellow equilibrium line $y=4$ will all tend to $y=4$ as $t\to\infty$. Thus, at any time $t$, if the $y$-values ever deviate slightly from $y=4$, we can be sure that the $y$-values will eventually tend toward $4$ again, which tells us that the $y$-values behave in a stable manner around $y=4$. 

I like to think of a stable equilibrium line as a tight rope where the tight rope walker will \textit{always} find their balance even if they wobble slightly off-balance while walking the line! Another analogy is a person moving in orbit around a planet like Earth; a stable equilibrium orbit would mean that even if you moved slightly toward or away from Earth, you would always fall back into the same stable orbit!
  
  In contract, the equilibrium solution $y(t)=-1$ is \fbox{unstable}, since solution curves above and below the red equilibrium line $y=4$ will all move \textit{away from} $y=-1$ as $t\to\infty$. Thus, at any time $t$, if the $y$-values ever deviate slightly from $y=-1$, they will quickly fly away from $-1$ and never return to $-1$ again, which tells us that the $y$-values behave in an unstable manner around $y=-1$.  
  
In the tight rope analogy, an unstable equilibrium line would mean that the tight rope walker will immediately lose their balance and fall off the rope if they wobble even the tiniest bit while walking the line! In the orbit analogy, an unstable equilibrium line would mean that if you moved slightly out of orbit either toward or away from Earth, you would either plummet toward Earth or fly off into space!

  \end{solution}

  \item
  \begin{problem}[5in]
  $\displaystyle \frac{dy}{dt} = (y + 1) (y - 2)^2$
  \end{problem}

  \begin{solution}
  We start by finding the equilibrium solutions, where $\frac{dy}{dt}
  = (y + 1)(y - 2)^2 $ is 0.  This happens when $y = -1$ or $y = 2$,
  so $y(t) = -1$ and $y(t) = 2$ are equilibrium solutions.

  Now, we need to figure out whether $\frac{dy}{dt}$ is positive or
  negative when $y < -1$, $-1 < y < 2$, and when $2 < y$.  To do this,
  we can test one point in each region, and we find:
  \begin{itemize}[topsep=0pt,itemsep=0pt]
  \item When $y < -1$, $\frac{dy}{dt} < 0$.  (For instance, plug in $y
  = -2$: $\frac{dy}{dt} = (y + 1) (y - 2)^2$ is equal to $(-1)
  (-4)^2 = -16$ here.)
  \item When $-1 < y < 2$, $\frac{dy}{dt} > 0$.
  \item When $2 < y$, $\frac{dy}{dt} > 0$.
\end{itemize}
We can put this information in a diagram:
\begin{center}
\includegraphics[width=2in]{ws/images/qa3-sketch}
\end{center}
The arrows indicate the general direction (up or down) of the
solution curves in each strip.  Now, it is easy to draw a
representative solution in each strip:
\begin{center}
\includegraphics[width=3in]{ws/images/qa3}
\end{center}

In this case, the equilibirum solution $y(t)=-1$ is \fbox{unstable}, since solution curves above and below the red equilibrium line move \textit{away from} the red line as $t\to\infty$. 

The equilibrium solution $y(t)=2$ is also \fbox{unstable}, because even though solution curves below the green line tend toward the green line as $t\to\infty$, solution curves above the green line move \textit{away from} the green line as $t\to\infty$. Thus, if our $y$-values ever deviated from $y=2$ in the positive direction, they would never return to $y=2$, which should feel like instability. (In our tight rope walker analogy, this would represent a tight rope walker who would always find their balance again if they wobbled toward one side of the rope, but would immediately lose their balance and fall off if they wobbled toward the other side! In the orbit analogy, this would mean that if you moved slightly toward Earth, you would still fall back into the same orbit, but if you moved slightly away from Earth, you'd fly off into space! In either case, I doubt you would feel very stable on the tight rope or in the orbit around Earth!)

Note that what we get from the qualitative analysis is just the
picture we'd get from drawing the slope field and following it:
\begin{center}
\includegraphics[width=3in]{ws/images/qa3-slope-field}
\end{center}
The advantage of using qualitative analysis is simply that we save a
lot of time.  Essentially, when we do qualitative analysis, all we
do is look at whether the slopes are positive, negative, or 0; we
can then infer the entire picture from just that information.
\end{solution}

\end{enumerate}    
\probonly{\end{multicols}}

\pspace{\newpage}



\item 
%\note{There are a few ways to have students engage in this problem. One way is to have students pair up (or get into small groups) and split the graphs between them. Each person can work on finding a differential equation for their graphs, as well as act as a checker for their partner's work. Another way is to put students into groups, and assign a graph to each group. Then have each group present their graph and how they arrived at their differential equation.}

For each picture, think of an autonomous first-order differential
equation whose solution curves look like those in the picture.
\probonly{\begin{multicols}{2}}
\begin{enumerate}
  \item
  \begin{problem}[2in]
  \raisebox{1ex-\height}{\includegraphics[width=2in]{ws/images/qa1-1}}
  \end{problem}

  \begin{solution}
  In this type of problem, there is no one correct answer.

  However, we can see from the picture that there is only one
  equilibrium solution, $y = 2$, and we should have
  \begin{itemize}[topsep=0pt,itemsep=0pt]
  \item $\frac{dy}{dt} > 0$ when $y > 2$.
  \item $\frac{dy}{dt} < 0$ when $y < 2$.
\end{itemize}
Or, if you prefer to look at a diagram:
\begin{center}
\includegraphics[width=2in]{ws/images/qa1-1-sketch}
\end{center}
So, one possibility is $\frac{dy}{dt} = y - 2$.

Getting a reasonable differential equation in such a problem is very
much a process of trial and error: write down a differential
equation, and see if it matches the given picture.  If not, see if
multiplying by $-1$ helps or squaring one or more factors helps.
\end{solution}

\item
\begin{problem}[2.3in]
\raisebox{1ex-\height}{\includegraphics[width=2in]{ws/images/qa1-2}}
\end{problem}

\begin{solution}
One possibility is $\frac{dy}{dt} = (y - 2)^2$.
\end{solution}

\item
\begin{problem}[2in]
\raisebox{1ex-\height}{\includegraphics[width=2in]{ws/images/qa1-3}}
\end{problem}

\begin{solution}
One possibility is $\frac{dy}{dt} = -(y - 2)^2$.
\end{solution}

\item
\begin{problem}[2in]
\raisebox{1ex-\height}{\includegraphics[width=2in]{ws/images/qa1-4}}
\end{problem}

\begin{solution}
One possibility is $\frac{dy}{dt} = -(y - 10)(y - 30)$.  (Note that
the equilibrium solutions are $y = 10$ and $y = 30$, so
$\frac{dy}{dt}$ should factor with linear factors of $y - 10$ and $y
- 30$.)
\end{solution}
\pspace{\newpage}
\item
\begin{problem}[1in]
\raisebox{1ex-\height}{\includegraphics[width=2in]{ws/images/qa1-5}}
\end{problem}

\begin{solution}
One possibility is $\frac{dy}{dt} = - (y - 5) (y - 10) (y + 5)$.
\end{solution}

 \item
 \begin{problem}[1in]
 \raisebox{1ex-\height}{\includegraphics[width=2in]{ws/images/qa1-6}}
 \end{problem}

 \begin{solution}
 One possibility is $\frac{dy}{dt} = (y - 5) (y - 10) (y + 5)$.
 \end{solution}

\item
\begin{problem}[1in]
\raisebox{1ex-\height}{\includegraphics[width=2in]{ws/images/qa1-7}}
\end{problem}

\begin{solution}
One possibility is $\frac{dy}{dt} = (y - 2)^2 (y + 1)^2$.
\end{solution}

 \item
 \begin{problem}[1in]
 \raisebox{1ex-\height}{\includegraphics[width=2in]{ws/images/qa1-8}}
 \end{problem}

 \begin{solution}
 One possibility is $\frac{dy}{dt} = (y - 1) (y + 2)^2$.
 \end{solution}

\end{enumerate}
\probonly{\end{multicols}}

\end{enumerate}

\end{document}